% Part: computability
% Chapter: recursive-functions
% Section: introduction

\documentclass[../../../include/open-logic-section]{subfiles}

\begin{document}

\olfileid[pt]{cmp}{rec}{int}
\olsection{Introdução}

A fim de desenvolver uma teoria matemática da computabilidade, é preciso,
antes de tudo, desenvolver um \emph{modelo} de computabilidade. Hoje pensamos
na computabilidade como o tipo de coisa que os computadores fazem, e os
computadores trabalham com símbolos. Contudo, no início do desenvolvimento
das teorias da computabilidade, o exemplo paradigmático de computação era a
computação \emph{numérica}. Os matemáticos sempre se interessaram por funções
teórico-numéricas, isto é, funções $f\colon \Nat^n \to \Nat$ que podem ser
calculadas. Portanto, não surpreende que, no início da teoria da
computabilidade, fossem essas as funções estudadas. Os exemplos mais
familiares de funções numéricas computáveis, como a adição, a multiplicação e
a exponenciação (de números naturais), compartilham uma característica
interessante: podem ser definidos \emph{recursivamente}. Assim, é bastante
natural tentar formular uma definição geral de \emph{função computável} com
base em definições recursivas. Entre as muitas maneiras possíveis de definir
recursivamente funções teórico-numéricas, um padrão de definição
particularmente simples torna-se central aqui: a chamada \emph{recursão
primitiva}.

Além de funções computáveis, podemos nos interessar por conjuntos e relações
computáveis. Um conjunto é computável se podemos calcular a resposta à
pergunta de se um dado número é ou não !!a{element} do conjunto; uma relação
é computável se podemos calcular se uma tupla
$\tuple{n_1, \dots, n_k}$ é ou não !!a{element} da relação. Considerando a
\emph{função característica} de um conjunto ou de uma relação, a discussão
de conjuntos e relações computáveis pode ser abrangida pela discussão de
funções computáveis. Assim, também podemos definir relações recursivas
primitivas; por exemplo, a relação ``$n$ divide $m$ sem resto'' é uma relação
recursiva primitiva.

As funções recursivas primitivas---aquelas que podem ser definidas usando
apenas recursão primitiva---não são, contudo, as únicas funções
teórico-numéricas computáveis. Muitas generalizações da recursão primitiva
foram consideradas, mas o modo adicional mais poderoso e amplamente aceito
de calcular funções é a busca ilimitada. Isso leva à definição de
\emph{funções recursivas parciais} e a uma definição relacionada, a de
\emph{funções recursivas gerais}. As funções recursivas gerais são
computáveis e totais, e a definição caracteriza exatamente as funções
recursivas parciais que por acaso são totais. As funções recursivas podem
simular todos os outros modelos de computação (máquinas de Turing, cálculo
lambda etc.) e, portanto, representam um dos muitos modelos de computação
aceitos.


\end{document}
